\documentclass[12pt]{article}
\usepackage[utf8]{inputenc}
\usepackage{amsmath}
\usepackage{graphicx}
\usepackage{hyperref}
\usepackage{geometry}
\geometry{margin=1in}
\usepackage{titlesec}

% Formatting for section headers
\titleformat{\section}{\bfseries\Large}{\thesection}{1em}{}
\titleformat{\subsection}{\bfseries\large}{\thesubsection}{1em}{}
\titleformat{\subsubsection}{\bfseries\normalsize}{\thesubsubsection}{1em}{}

\title{\textbf{Deconvolution of Clinical Variance in CAR-T Cell Pharmacology and Response}}
\author{}
\date{}

\begin{document}

\maketitle

\section*{Introduction}
This paper investigates the variability in outcomes for CAR-T cell therapies. These therapies have revolutionized cancer treatment but exhibit significant variability in efficacy and toxicity across patients. The study aims to:
\begin{itemize}
    \item Understand the mechanisms driving this variability.
    \item Develop predictive models to guide CAR-T product optimization.
\end{itemize}

\section*{Methodology}

\subsection*{Mathematical Model}
\begin{itemize}
    \item A toggle switch model regulates CAR-T cells transitioning among memory (TM), effector (TE), and exhausted (TX) states based on tumor antigen engagement.
    \item Parameters were estimated using data from clinical trials, such as Kymriah-treated chronic lymphocytic leukemia (CLL) patients.
    \item The model was validated across patient response categories: complete responders (CR), partial responders (PR), and non-responders (NR).
\end{itemize}

\subsection*{Transcriptomic Analysis}
\begin{itemize}
    \item Bulk RNA-seq data analyzed pre-infusion CAR-T products for differential gene expression between CR and NR groups.
    \item Single-cell RNA-seq (scRNA-seq) provided finer insights into heterogeneity, exhaustion, and memory-related gene signatures.
\end{itemize}

\subsection*{Machine Learning Workflow}
\begin{itemize}
    \item Logistic regression classifiers were trained on transcriptomic data to predict CR vs. NR outcomes.
    \item The model's predictive accuracy was compared to immunophenotyping data (e.g., T cell subpopulations like memory and exhausted cells).
\end{itemize}

\section*{Results}

\subsection*{Key Determinants of Response}
\begin{itemize}
    \item Memory T cell turnover and effector cell cytotoxicity were critical to CAR-T expansion and tumor clearance.
    \item CRs exhibited higher memory cell function and lower exhaustion markers (e.g., LAG3, PD1) compared to NRs.
\end{itemize}

\subsection*{Gene Signatures for Response Prediction}
\begin{itemize}
    \item CR products were enriched in IL2RB, IL7R, and other memory T cell-related pathways.
    \item NR products showed heightened DNA damage and p53 signaling pathways, correlating with reduced proliferative and cytotoxic capacities.
\end{itemize}

\subsection*{Transcriptomic Classifier Performance}
\begin{itemize}
    \item The transcriptomic classifier achieved 90\% accuracy for predicting outcomes, outperforming immunophenotyping-based methods (60--83\% accuracy).
    \item A CAR-T response scorecard highlighted key pathways shared across datasets.
\end{itemize}

\subsection*{Simulating Dose-Response Variability}
\begin{itemize}
    \item Simulations showed that the ratio of CAR-T expansion (Cmax) to tumor burden (B0) predicts response durability.
    \item Applied the model to other therapies, like Abecma for multiple myeloma, validating its generalizability.
\end{itemize}

\section*{Discussion}

\subsection*{Clinical Implications}
\begin{itemize}
    \item Pre-infusion transcriptomics can guide CAR-T optimization and predict patient outcomes.
    \item Transitioning to allogeneic CAR-T therapies may reduce variability caused by autologous starting materials.
\end{itemize}

\subsection*{Future Directions}
\begin{itemize}
    \item Expand datasets to validate findings across more CAR-T products and tumor types.
    \item Incorporate host-intrinsic factors, like immune microenvironment and tumor signaling, into models.
    \item Use these insights for personalized dose optimization.
\end{itemize}

\section*{Figures and Highlights}
\begin{itemize}
    \item \textbf{Model Outputs:} PCA plots show parameter-driven separation of CR, PR, NR populations.
    \item \textbf{Gene Signature Analysis:} Heatmaps of pathways differentiating CR and NR products.
    \item \textbf{Simulations:} Dose-response curves demonstrating relationships between Cmax, tumor burden, and clinical outcomes.
\end{itemize}

\section*{Conclusion}
This study integrates mathematical modeling, transcriptomics, and machine learning to elucidate and predict clinical variability in CAR-T therapies. It identifies product-intrinsic factors, such as memory cell function and exhaustion, as critical determinants of efficacy, paving the way for reliable CAR-T designs and individualized treatment strategies.

\end{document}
